\documentclass[aps,prb,onecolumn,nofootinbib,superscriptaddress]{revtex4-2}
\usepackage{amsmath,amssymb,amsthm,mathtools,bm}
\usepackage{booktabs}
\usepackage[colorlinks=true,linkcolor=blue,citecolor=blue,urlcolor=blue]{hyperref}
\usepackage{microtype}

\newcommand{\I}{\mathbf{1}}
\newcommand{\Tr}{\operatorname{Tr}}
\newcommand{\dd}{\mathrm{d}}
\newcommand{\calC}{\mathcal{C}}
\newcommand{\Dlam}{\Delta_{\lambda}}
\newcommand{\Tp}{T_{\rm p}}
\newcommand{\SigmaLL}{\Sigma_{LL}}
\newcommand{\SigmaLR}{\Sigma_{LR}}
\newcommand{\pathref}[1]{\path{#1}}

\begin{document}

\title{Supervisor 04 Synthesis: Transfer Stability, Choi Gaps, and Lyapunov Evidence}
\author{Supervisor 04}
\date{\today}

\begin{abstract}
This note synthesizes the assigned worker reports on Lyapunov spectra,
regularized GPU Choi transfer matrices, and CPU particle-Choi transfer gaps,
and checks those claims against the non-excluded transfer/stability source
areas of the repository.  The repository supports a robust qualitative
conclusion: the uniform topological geometry is strongly contracting in the
measured transfer diagnostics, whereas the domain-wall geometry exhibits a
near-marginal slow channel localized to, or at least naturally associated with,
the domain-wall sector.  The strongest numerical evidence is the finite-time
single-particle Lyapunov campaign, where the domain-wall gap is suppressed by
one to two orders of magnitude relative to the uniform geometry and decreases
over \(N_y=30,40,50\).  The Choi-transfer evidence is complementary rather
than identical: the GPU regularized Choi chart gives an early-cycle finite-sector
diagnostic but censors before final time, while the CPU Choi pipeline records
stable final Choi spectra and eigenvectors for slab-restricted domain-wall runs
with a separate endpoint-saturation caveat.  The Boss should preserve the
qualitative convergence of diagnostics and scrutinize any claim of an already
resolved thermodynamic scaling law or equality between the Lyapunov and Choi
operators.
\end{abstract}

\maketitle
\setcounter{tocdepth}{2}
\tableofcontents

\section{Source Scope}

I read \path{REPO_POLICY.md} before inspecting the repository.  I did not edit
files, execute notebooks, or load large binary arrays.  I excluded every path
containing \texttt{haining}, \texttt{haoyu}, or \texttt{erroneous},
case-insensitively.  I read only the assigned worker notes
\path{repo_synthesis/main_pass/worker_notes/worker_12_lyapunov_spectra.tex},
\path{repo_synthesis/main_pass/worker_notes/worker_13_regularized_choi_transfer.tex},
and
\path{repo_synthesis/main_pass/worker_notes/worker_14_cpu_choi_transfer_gap.tex};
I did not read other worker notes, supervisor notes, or boss notes.

The independent audit covered the requested transfer/stability areas:
\path{colab_lyapunov/},
\path{colab_regularized_choi_transfer_matrix/},
\path{choi_covariance_cpu/}, and
\path{lyapunov_analysis_v2/}.  I also inspected non-excluded transfer/gap
scripts and documents surfaced by a repository search, including
\path{markov_transfer_operators/markov_transfer_operators_note.tex} and
\path{scripts/smoke_choi_dense_vs_rank1.py}.  Binary spectra and eigenvector
payloads such as \texttt{.npz} files were not opened.

\section{Synthesis of Worker Claims}

\subsection{Worker 12: Lyapunov Spectra}

Worker 12's account is well supported by the audited sources.  The Lyapunov
campaign measures a Benettin/QR spectrum for a frame
\(V\in\mathbb{C}^{2N_xN_y\times n_{\rm vec}}\) living in the top-layer
single-particle orbital basis, not the full real tangent space of Hermitian
covariance perturbations.  The implemented diagnostic is the final-time
finite-cycle gap
\begin{equation}
  \Dlam(t)=\min_i |\lambda_i(t)|,
\end{equation}
where the logarithmic singular accumulation is normalized per Markov cycle.

The key numerical claim is that \(DW=0\) is \(O(1)\)-gapped, while \(DW=1\) is
near-marginal.  The writeup reports, for perfect correction,
\[
  \Dlam_{DW0}^{\rm PC}(N_y=30,40,50)
  \simeq (0.722,0.748,0.768),
  \qquad
  \Dlam_{DW1}^{\rm PC}
  \simeq (0.027,0.024,0.020),
\]
and for post-selection
\[
  \Dlam_{DW0}^{\rm PS}\simeq(0.945,0.955,0.952),
  \qquad
  \Dlam_{DW1}^{\rm PS}\simeq(0.021,0.0076,0.0078).
\]
The same source reports that \(DW=0\) plateaus within roughly ten cycles, while
\(DW=1\) does not clearly plateau within \(T=N_y\).  This is the central
stability contrast.

Worker 12 also reports a physical localization/correlation story.  The Chern
marker remains quantized for \(DW=0\), is sub-integer for \(DW=1\), and
sample-by-sample smaller \(\calC\) correlates with smaller \(\Dlam\).  The
minimum-\(|\lambda|\) vector is described as concentrated near the domain-wall
columns \(x=5,15\) and extended along \(y\).  Because I did not load the binary
eigenvector arrays, the spatial statement remains accepted as text-level
writeup evidence rather than independently recomputed evidence.

\subsection{Worker 13: Regularized GPU Choi Transfer}

Worker 13's interpretation is also supported, with one important limitation.
The regularized Choi construction is an operator-history chart for a realized
Gaussian circuit operator, not the physical covariance trajectory itself.  The
audited document distinguishes the two-leg Choi covariance
\[
  \Sigma =
  \begin{pmatrix}
  \Sigma_{LL} & \Sigma_{LR}\\
  \Sigma_{RL} & \Sigma_{RR}
  \end{pmatrix}
\]
from the production state covariance \(G\).  The local singular Choi
contraction is regulated by replacing the \(Q\)-sector null block with an
\(\epsilon Q\) inverse and taking the Schur-complement limit.  The resulting
rank-one formula depends on the nonsingular denominator
\[
  d=\langle\chi|\Sigma_{RR}|\chi\rangle-\eta_1.
\]
The GPU implementation initializes
\(\Sigma_{LL}=0\), \(\Sigma_{LR}=\I\), \(\Sigma_{RR}=0\), tracks the minimum
\(|d|\), and records or censors small-denominator/unstable branches.

The finite-sector particle transfer diagnostic is
\begin{equation}
  \Tp=-(\I+\SigmaLL)^{-1}\SigmaLR,
  \qquad
  \Delta_{\rm p}(t)=\min_j\left|\frac{\log s_j(\Tp)}{t}\right|.
\end{equation}
In the locally present GPU summaries and CSV aggregation, the \(DW=1\) slab
mean gaps during the stable window are roughly \(10^{-2}\) to
\(2\times10^{-2}\), while the homogeneous \(DW=0\) full control has much larger
transient gaps after the first few cycles.  However, all locally available GPU
Choi samples are censored before the final cycle: the stable fraction reaches
zero near cycles \(15\)--\(16\).  The final \(T=2N_y\) GPU Choi gap is therefore
not a resolved late-time physical gap; it is intentionally absent after Choi
chart failure.

\subsection{Worker 14: CPU Choi Transfer Gap}

Worker 14 correctly identifies the CPU pipeline as the strongest Choi-side
production evidence.  The CPU script uses
\texttt{classA\_U1FGTN.run\_markov\_circuit} and records the same particle
definition
\[
  \Tp=-(\I+\SigmaLL)^{-1}\SigmaLR,
  \qquad
  \Delta_{\rm p}(C)=\min_j\left|\frac{\log s_j(\Tp)}{C}\right|.
\]
Operationally, it diagonalizes \(\SigmaLL\) exactly and uses
\[
  \lambda_{\rm p}(a;C)=
  \frac{\log(1-a)-\log(1+a)}{2C}
\]
for finite \(\SigmaLL\)-eigenvalues \(a\in(-1,1)\), while treating endpoints
as particle zeros or poles.

The \(N_x=20\), \(N_y=20,30,40\), \(DW=1\), slab-restricted CPU campaign with
\(\alpha_1\in\{1,3\}\), \(\alpha_2=30\), \(S=10\), and \(C=2N_y\) is complete
in the metadata.  Its run index records six completed configurations, ten
stable final samples in each, rank-one physical covariance updates, and no
dense fallback.  The newer metadata also records that final finite eigenstates
are saved and ordered by increasing \(|\log s|\).

The caveat is endpoint saturation.  The CPU observer policy explicitly accepts
cycles with fewer than the requested number of finite transfer exponents and
pads missing near-gap slots with NaNs and \(-1\) indices.  Thus
``stable final samples'' means the Choi chart survived, not necessarily that
each run has three resolved finite near-gap modes.  Older/executed plotting
notebooks also show sentinel behavior in some \(\alpha_1=3\) runs, so endpoint
collapse must be separated from a finite nonzero gap.

\section{Independent Repo-Audit Adjustments}

First, the transfer word is used in several inequivalent senses.  The
\path{markov_transfer_operators} note formulates a Markov kernel on covariance
matrices and a full covariance tangent cocycle
\[
  \delta G_{t+1}=T_{\alpha_t}(G_t)^\dagger\delta G_t T_{\alpha_t}(G_t).
\]
By contrast, the Colab Lyapunov implementation explicitly propagates
single-particle frame columns \(v\in\mathbb{C}^{2N_xN_y}\).  The Boss should
not collapse these into one operator.  The implemented Lyapunov spectrum is a
single-particle transfer-frame diagnostic; it is not the full Hermitian
\(\delta G\) Lyapunov spectrum.

Second, the Choi-transfer evidence is not uniform across GPU and CPU.  The GPU
regularized Choi campaign is valuable as a controlled early-cycle diagnostic,
but its local summaries censor all samples before final time because Choi
entries leave the accepted numerical chart.  The CPU campaign, by contrast,
does reach final cycles with stable Choi charts, but it uses an endpoint policy
that can leave fewer finite modes than requested.  These are different
failure/saturation modes and should not be merged.

Third, there is minor version drift in the Choi metadata.  The current GPU
helper advertises a resolvent/rank-one version string, while some run summaries
record an earlier basis-independent formula label.  The actual audited code and
derivation are consistent in structure--rank-one denominator, resolvent action,
finite-sector transfer from \(\SigmaLL\)--but exact provenance should quote the
run metadata rather than silently updating old runs to the newest helper name.

Fourth, the older \path{lyapunov_analysis_v2} Majorana Choi notebook should be
treat as historical context only.  It loads cached Choi covariances, converts
complex blocks to Majorana blocks, and often uses a pseudoinverse with very low
reported numerical rank, for example rank \(8/1280\) in the inspected notebook
outputs.  It is weaker evidence than the later CPU exact-\(\SigmaLL\) pipeline.

Fifth, the Lyapunov numerical trend is strong but finite-time.  The \(DW=1\)
gap decreases over three circumferences and is far below \(DW=0\), but the
writeup itself states that \(T=N_y\) does not establish a plateau or distinguish
\(1/N_y\), another power, or logarithmic closing.

\section{Key Evidence Table}

\begin{table*}[t]
\caption{\label{tab:evidence}Key text, metadata, and source-code evidence used in this synthesis.}
\begin{ruledtabular}
\begin{tabular}{p{0.30\linewidth}p{0.16\linewidth}p{0.46\linewidth}}
Source & Lines & Evidence \\
\colrule
\pathref{REPO_POLICY.md} & 1--37 & Canonical CPU/GPU dynamics entry points and repository constraints for GPU and CPU circuit simulations. \\
\pathref{colab_lyapunov/docs/lyapunov_analysis.tex} & 142--172 & Defines the propagated Lyapunov frame as a single-particle Hilbert-space object, not the full Hermitian covariance tangent space. \\
\pathref{colab_lyapunov/docs/lyapunov_analysis.tex} & 183--221 & QR accumulation, cycle normalization, \(\Dlam=\min_i|\lambda_i|\), and final min-gap vector storage. \\
\pathref{colab_lyapunov/docs/lyapunov_analysis.tex} & 223--230 & Campaign parameters: \(N_x=20\), \(N_y=30,40,50\), \(T=N_y\), \(S=25\) for perfect correction, \(S=1\) for post-selection, canonical GPU entry point. \\
\pathref{colab_lyapunov/docs/lyapunov_analysis.tex} & 258--301 & \(DW=0\) spectra are strictly negative and plateau quickly; \(DW=1\) develops near-zero modes and does not clearly plateau. \\
\pathref{colab_lyapunov/docs/lyapunov_analysis.tex} & 306--355 & Final Lyapunov gap table and unresolved scaling forms. \\
\pathref{colab_lyapunov/docs/lyapunov_analysis.tex} & 383--403, 464--471 & \(DW=1\) Chern marker degradation and sample-level correlation with smaller \(\Dlam\). \\
\pathref{colab_lyapunov/docs/lyapunov_analysis.tex} & 473--502 & Text-level spatial interpretation of the neutral vector as concentrated near domain-wall columns. \\
\pathref{colab_lyapunov/src/lyapunov_observables_gpu.py} & 13--15, 206--215, 267--350 & Helper version, canonical GPU entry point, stored spectra/Chern/charge arrays, min-abs vector payload, and top-layer complex orbital basis description. \\
\pathref{colab_lyapunov/gpu_data/.../run_index.csv} & 1--13 & Confirms the 12 Lyapunov configurations, actual sample counts, cycles, and \(n_{\rm vec}=2N_xN_y\). \\
\pathref{colab_regularized_choi_transfer_matrix/docs/regularized_choi_covariance_contractions.tex} & 63--103 & Separates cumulative operator Choi covariance from physical state covariance. \\
\pathref{colab_regularized_choi_transfer_matrix/docs/regularized_choi_covariance_contractions.tex} & 164--246 & Shows the formally singular two-leg contraction and regulator construction. \\
\pathref{colab_regularized_choi_transfer_matrix/src/classA_U1FGTN_gpu.py} & 754--774, 840--916 & Initializes Choi blocks, tracks \(|d|\), records denominator failures, and applies the rank-one/resolvent update. \\
\pathref{colab_regularized_choi_transfer_matrix/src/choi_transfer_observables_gpu.py} & 11--18, 43--63 & Records \(\Tp\), gap definition, endpoint treatment, and exact final-spectrum extraction convention. \\
\pathref{colab_regularized_choi_transfer_matrix/analysis_outputs/gap_vs_cycle/gap_vs_cycle_summary.csv} & 62--77, 122--136 & \(DW=1\) GPU Choi slab gaps are \(O(10^{-2})\) while stable samples disappear by cycles \(15\)--\(16\). \\
\pathref{colab_regularized_choi_transfer_matrix/gpu_data/.../N20x30_DW1.../run_summary.json} & 1426--1496, 1716--1779 & GPU Choi failure mode, formula metadata, first failure cycles, stable-fraction collapse, and zero final stable samples. \\
\pathref{colab_regularized_choi_transfer_matrix/gpu_data/.../N20x40_DW1.../run_summary.json} & 1646--1716, 1936--2019 & Same censoring pattern for \(N_y=40\). \\
\pathref{choi_covariance_cpu/run_particle_choi_transfer_gap_cpu.py} & 61--72 & CPU class import, canonical CPU entry point, \(\Tp\), gap, exact-\(\SigmaLL\) extraction, and rooted singular convention. \\
\pathref{choi_covariance_cpu/run_particle_choi_transfer_gap_cpu.py} & 241--306 & Exact Hermitianization/eigensolve of \(\SigmaLL\), endpoint handling, exponent formula, and near-gap selection. \\
\pathref{choi_covariance_cpu/run_particle_choi_transfer_gap_cpu.py} & 560--610 & Censoring, exact final storage, final finite eigenstate storage, and stability mask. \\
\pathref{choi_covariance_cpu/run_particle_choi_transfer_gap_cpu.py} & 844--904 & Calls \texttt{classA\_U1FGTN.run\_markov\_circuit} with \texttt{track\_choi=True} and the Choi observer. \\
\pathref{choi_covariance_cpu/cpu_data/.../Nx20...allfiniteeig_v4/run_index.csv} & 1--7 & Six completed \(N_x=20\), \(N_y=20,30,40\), \(\alpha_1=1,3\) slab \(DW=1\) CPU runs with ten stable final samples each. \\
\pathref{choi_covariance_cpu/cpu_data/.../Nx20...allfiniteeig_v4/campaign_manifest.json} & 1--17, 131--141 & Campaign geometry, canonical CPU entry point, final finite eigenstate saving, and partial finite-mode policy. \\
\pathref{choi_covariance_cpu/cpu_data/.../N20x40...a1-1.../run_summary.json} & 1200--1250, 1331--1357 & Representative final-stable CPU Choi metadata: no full matrices saved, exact extraction, no observer failures, stable final samples, and transfer object. \\
\pathref{choi_covariance_cpu/plot_gap_vs_cycle_nan_sentinel_executed.ipynb} & 141--147, 824--837 & Existing notebook outputs showing sentinel behavior in older \(\alpha_1=3\) curves despite stable final bookkeeping. \\
\pathref{choi_covariance_cpu/plot_nx16_alpha1_sweep_gap_vs_alpha.ipynb} & 57--64, 104--144 & Alpha-sweep plotting code treats nonfinite final gaps through an explicit sentinel and tracks finite final sample counts. \\
\pathref{choi_covariance_cpu/plot_nx16_alpha1_sweep_low_lying_modes.ipynb} & 87--103, 112--154 & Low-lying-mode plotting workflow loads final finite eigenstates and constructs spatial density maps, but this synthesis did not load those arrays. \\
\pathref{lyapunov_analysis_v2/majorana_choi_transfer_spectrum.ipynb} & 7--20, 165--201, 306--330 & Historical Majorana transfer notebook, cached Choi inputs, pseudoinverse fallback, and low-rank/ill-conditioned diagnostics. \\
\pathref{markov_transfer_operators/markov_transfer_operators_note.tex} & 36--48, 302--430 & Conceptual Markov-kernel and full covariance tangent-cocycle framework; useful for terminology but not identical to the implemented single-particle Lyapunov frame. \\
\pathref{scripts/smoke_choi_dense_vs_rank1.py} & 1--6, 158--176, 191--220, 253--331 & Small-system dense-versus-rank-one Choi health check and survival comparison script. \\
\end{tabular}
\end{ruledtabular}
\end{table*}

\section{Contradictions and Caveats}

There is no hard contradiction among the three worker notes, but there are
several places where a synthesis could overstate the evidence.

The Lyapunov and Choi gaps are not the same spectrum.  The Lyapunov campaign
measures a finite-time single-particle transfer-frame spectrum.  The Choi
campaigns measure finite-sector particle-transfer singular exponents extracted
from a cumulative Choi covariance.  The Markov-transfer note describes the
full covariance tangent cocycle.  These three objects are related in physical
motivation, but they are not interchangeable.

The GPU Choi results should not be described as final-cycle stable transfer
gaps.  Their physical \(G\) evolution continues, but the exact Choi observable
is censored after numerical chart failure.  The stable window is still useful:
it shows small \(DW=1\) slab gaps relative to the homogeneous control.  It does
not establish a late-time \(T=2N_y\) Choi gap.

The CPU Choi summaries are stronger at final time but carry an endpoint caveat.
Stable final samples mean the Choi chart survived; they do not by themselves
guarantee a fixed number of finite near-gap modes.  The Boss should preserve
the metadata distinction between chart stability, finite-mode count, NaN
padding, and sentinel plotting.

The spatial eigenvector/eigenstate story is promising but not independently
recomputed here.  The Lyapunov writeup reports domain-wall localization, and
the CPU notebooks define density-map workflows from saved Choi eigenvectors.
This synthesis did not load those binary eigenvector arrays.

The thermodynamic scaling law is unresolved.  Three \(N_y\) values and
finite-time runs with \(T=N_y\) support a qualitative closing tendency in the
domain-wall sector, but not a reliable exponent or a distinction between power
law and logarithmic closing.

\section{Confidence}

Confidence is high that the repository contains a coherent qualitative
transfer-stability story: \(DW=0\) is strongly contracting in the measured
Lyapunov diagnostic, while \(DW=1\) supports a much slower near-zero channel
linked to domain-wall physics.

Confidence is high for the implementation-level definitions of the diagnostics:
the audited source files explicitly record the canonical dynamics entry points,
the Lyapunov frame normalization, the Choi finite-sector transfer definition,
the exact-\(\SigmaLL\) extraction, and the censoring/endpoint policies.

Confidence is moderate for the physical identification of all near-gap
diagnostics with the same chiral domain-wall mode.  The convergence of evidence
is persuasive, but the operators are not identical and the spatial evidence was
not recomputed from binary arrays.

Confidence is low for any precise thermodynamic scaling law or universal
critical exponent at this stage.

\section{What the Boss Should Preserve or Scrutinize}

The Boss should preserve the central qualitative statement: independent
Lyapunov and Choi-transfer diagnostics point to a near-marginal domain-wall
sector, while the uniform topological geometry remains strongly stable.

The Boss should preserve the operator distinctions.  Use phrases such as
``single-particle Lyapunov transfer-frame gap'' and ``finite-sector Choi
particle-transfer gap'' rather than a single generic ``transfer gap''.

The Boss should preserve the strongest numbers from the Lyapunov table, since
they give the cleanest \(DW=0\) versus \(DW=1\) contrast over
\(N_y=30,40,50\).

The Boss should scrutinize any use of GPU Choi final-cycle gaps, because the
local GPU Choi summaries have zero stable final samples.

The Boss should scrutinize CPU Choi claims that ignore endpoint saturation,
NaN padding, or finite-mode counts.

The Boss should scrutinize any claim that the scaling is already known, or that
the Lyapunov and Choi spectra are mathematically identical.  The evidence
supports qualitative convergence of diagnostics, not equality of operators.

\section{Cross-Supervisor Addendum}

\subsection{Cross-note agreements}

All five supervisor notes agree on the core object hierarchy.  The physical
adaptive circuit evidence should be anchored to the canonical Markov drivers,
\texttt{classA\_U1FGTN.run\_markov\_circuit(...)} and
\texttt{classA\_U1FGTN\_gpu.run\_markov\_circuit(...)}, with perfect correction
understood as a Born-sampled trajectory ensemble with deterministic feedback,
not as postselection or an outcome-averaged transfer matrix.  This directly
supports the distinction made here between sampled Lyapunov products, Choi
operator-transfer diagnostics, and distributional transfer-operator language.

The target-state and trajectory notes reinforce the physical interpretation of
the transfer/stability evidence.  Supervisor 01 identifies the domain-wall
benchmark and OW target-state structure as chiral-wall physics, while Supervisor
03 finds trajectory-level entanglement slopes near \(1/3\) and strong
perfect-correction purification.  These are consistent with the Supervisor 04
claim that the near-zero Lyapunov/Choi modes are domain-wall-sector slow
directions rather than generic bulk instabilities.

Supervisor 05 agrees with the evidentiary hierarchy used here: structured
campaign manifests, run summaries, and run-index CSV files are the safest data
spine; legacy caches and hard-coded or notebook-only products are weaker.  This
supports preserving the Lyapunov campaign metadata and the CPU Choi run-index
evidence, while treating older Majorana Choi notebooks and raw cache artifacts
as historical context.

\subsection{Conflicts and corrections to Supervisor 04}

The main correction is provenance strength.  Supervisor 02 shows that a recorded
canonical-entry string does not by itself prove byte identity with the canonical
\texttt{src/fgtn} source: several Colab/OSG GPU copies are stale.  Therefore
the version-drift caveat in this note should be read more strongly.  Choi
formula labels, helper versions, and copied-source identity must be quoted from
the actual run metadata, and high-stakes claims should prefer outputs with
source-copy or checksum provenance.

A second correction is convention hygiene.  Supervisor 01 reports that
\(\alpha_1,\alpha_2\) are not used uniformly across source areas.  The
transfer note's statements about \(\alpha_1=1,\alpha_2=30\) are correct for the
audited Choi/Lyapunov campaigns, but Boss-level prose should translate these as
``topological slab mass'' and ``trivial flank mass'' rather than assuming a
global \(\alpha_1,\alpha_2\) convention.

A third correction is statistical framing.  Supervisor 03 and Supervisor 05 both
emphasize that postselection commonly has one actual trajectory even when a
larger sample count was requested.  The Supervisor 04 postselection caveat is
therefore not merely a Lyapunov-specific detail; it is a repository-wide sample
semantics issue.

No cross-note evidence weakens the main Supervisor 04 conclusion.  Instead, the
cross-share narrows its wording: the repository supports qualitative convergence
of domain-wall slow-mode diagnostics, not equality of Lyapunov and Choi
operators, not a final thermodynamic exponent, and not a postselection ensemble
comparison.

\subsection{Boss priorities}

\begin{enumerate}
\item Preserve the object distinctions: target-state benchmarks, sampled
perfect-correction trajectories, deterministic postselection controls, Choi
operator-transfer diagnostics, Lyapunov transfer-frame spectra, and
mean-channel/Lindblad objects are related but not interchangeable.

\item Use the structured manifest/run-summary layer as the primary evidence
spine, and require canonical-entry plus source-identity provenance for any
flagship production result.

\item Keep the domain-wall story qualitative but coherent: exact/OW target
states, trajectory entanglement, Lyapunov gaps, and Choi gaps all point toward a
slow chiral wall sector, while current data do not resolve a universal scaling
law.

\item Scrutinize postselection sample counts, GPU Choi final-cycle claims,
CPU Choi endpoint saturation, stale copied engine sources, and any use of
legacy cache files not tied to current manifests.

\item Write \(\alpha\)-parameter captions in physical language
(topological slab versus trivial flank) to avoid cross-folder convention errors.
\end{enumerate}

\end{document}
